% test-029.tex -- comandos \spProb y \spcomp
%
% Compilar:  lualatex-dev test-029.tex
% Validar:   show-pdf-tags --xml test-029.pdf | rnv-wrapp
%            verapdf --flavour ua2 --format text test-029.pdf
%
% Cada caso es un parrafo numerado, para poder referirse a el al
% escuchar el PDF. Los comentarios "doc:" son lecturas reales de
% NVDA con MathCAT (2026-10-09). Las entradas invalidas (que detienen la compilacion)
% no van aqui.
\DocumentMetadata
  {
    lang=es-CL,
    pdfversion=2.0,
    pdfstandard=ua-2,
    tagging=on,
    tagging-setup={math/setup={mathml-SE}},
  }
\documentclass{article}
\usepackage[spanish]{babel}
\usepackage{unicode-math}
\usepackage{spintent}
\usepackage{hyperref}
\hypersetup
  {
    pdfdisplaydoctitle = true,
    pdftitle           = {test-029: probabilidad de eventos},
  }
\pagestyle{empty}
\begin{document}

\section*{A. Un evento y dos eventos}

% doc: «probabilidad de A»
1. Un evento: $\spProb{A}$.

% doc: «probabilidad de A y B»
2. Intersección: $\spProb{A \cap B}$.

% doc: «probabilidad de A o B»
3. Unión: $\spProb{A \cup B}$.

% doc: «probabilidad de A dado B»
4. Condicional con mid: $\spProb{A \mid B}$.

% doc: «probabilidad de A dado B»
5. Condicional con barra: $\spProb{A | B}$.

\section*{B. Eventos con estructura}

% doc: «probabilidad de a subíndice 1 y a subíndice 2»
6. Con subíndices: $\spProb{A_{1} \cap A_{2}}$.

% doc: «probabilidad de a y complemento de b»
7. Con complemento: $\spProb{A \cap \spcomp{B}}$.

% doc: «probabilidad de complemento de a»
8. Complemento solo: $\spProb{\spcomp{A}}$.

% doc: «probabilidad de a y se abren paréntesis b o c se cierran paréntesis»
9. Con paréntesis: $\spProb{A \cap (B \cup C)}$.

% doc: «probabilidad de se abren paréntesis a o b se cierran paréntesis y c»
10. Con left y right: $\spProb{\left( A \cup B \right) \cap C}$.

\section*{C. Llaves}

% doc: «probabilidad del suceso»
11. Con read-arg: $\spProb[read-arg={probabilidad del suceso}]{A}$.

% doc: sin lectura
12. Con silent: $\spProb[silent=true]{A \cap B}$.

\section*{D. Dentro de fórmulas}

% doc: «probabilidad de a es igual a 1 de 5»
13. Con valor: $\spProb{A} = \spPval{1}{5}$.

% doc: «probabilidad de a es igual a casos favorables partido por casos posibles»
14. Regla de Laplace: $\spProb{A} = \dfrac{\spcase{f}}{\spcase{p}}$.

% doc: «probabilidad de a o b es igual a probabilidad de a más probabilidad de b menos probabilidad de a y b»
15. Unión de dos eventos: $\spProb{A \cup B} = \spProb{A} + \spProb{B} - \spProb{A \cap B}$.

% doc: «probabilidad de a y b es igual a probabilidad de a punto probabilidad de b dado a»
16. Producto: $\spProb{A \cap B} = \spProb{A} \cdot \spProb{B \mid A}$.

\section*{E. El comando spcomp}

% doc: «complemento de a»
17. Complemento: $\spcomp{A}$.

% doc: «complemento de a»
18. Con barra: $\spcomp[bar=true]{A}$.

% doc: «complemento de se abren paréntesis a o b se cierran paréntesis»
19. Complemento de una unión: $\spcomp{A \cup B}$.

% doc: «no ocurre A»
20. Con read-arg: $\spcomp[read-arg={no ocurre A}]{A}$.

% doc: sin lectura
21. Con silent: $\spcomp[silent=true]{A}$.

% doc: «probabilidad de a y complemento de b»
22. Con barra dentro de spProb: $\spProb{A \cap \spcomp[bar=true]{B}}$.

% doc: «probabilidad de complemento de a es igual a 1 menos probabilidad de a»
23. Regla del complemento: $\spProb{\spcomp{A}} = 1 - \spProb{A}$.

% doc: «probabilidad de complemento de se abren paréntesis a o b se cierran paréntesis es igual a probabilidad de complemento de a y complemento de b»
24. Complemento de la unión: $\spProb{\spcomp{A \cup B}} = \spProb{\spcomp{A} \cap \spcomp{B}}$.

% doc: «complemento de a o b»
25. Sin paréntesis: $\spcomp[paren=false]{A \cup B}$.

% doc: «complemento de a»
% limitación conocida: MathCAT omite los paréntesis alrededor de un solo símbolo
26. Con paréntesis forzados: $\spcomp[paren=true]{A}$.

% doc: «complemento de la expresión a o b»
27. Barra sobre una unión: $\spcomp[bar=true]{A \cup B}$.

% doc: «probabilidad de complemento de se abren paréntesis a y b se cierran paréntesis es igual a 1 menos probabilidad de a y b»
28. Complemento de una intersección: $\spProb{\spcomp{A \cap B}} = 1 - \spProb{A \cap B}$.

% doc: «probabilidad de complemento de la expresión a y b es igual a 1 menos probabilidad de a y b»
29. Barra sobre una intersección: $\spProb{\spcomp[bar=true]{A \cap B}} = 1 - \spProb{A \cap B}$.

\end{document}
